% Root draft; uninput pending independent literal audit.
\subsection{The seventeen-arrow three-object strong-component cases}

\begin{lemma}\label{C:small-corner-two-states-per-object}
Let $S$ be a finite category and let $N$ be a monoid of order
at most four. Form a collage with full source block $S$, sink
corner $N$, and compatible cross sets $F(i)$ satisfying
$1\leq|F(i)|\leq2$ for every object of $S$. Its classifying
space cannot have type $\mathbb CP^2$.
\end{lemma}
\begin{proof}
If the trivial right rational $N$-module is projective, the
actual orbit-collapse functor of
Proposition~\ref{C:finite-target-extreme-corner} is a rational
cohomology-ring equivalence to an ordinary categorical act cone
with at most two states at each object. Apply
Lemma~\ref{C:two-states-objectwise-act-cups}.

Otherwise the small-corner Rees argument in
Theorem~\ref{C:small-opposite-corner-three-states} says that the
minimal ideal of $N$ is a right-zero band of size two or three.
Its left actors are retractions onto the same subset $U(i)$ of
each $F(i)$. On a set of at most two elements these retractions
are identical: if the common image has one element they are
the same constant map, and if it has two they are identities.
Choose any minimal element $p$ and write $P_i$ for this common
map. Every $N$-actor fixes $U(i)$ pointwise. The identities
$np=p$ and $pn$ minimal therefore show
$nP_i=P_i=P_in$. The family $P_i$ also commutes with all
right $S$-act maps. It is consequently an actual biset retraction
onto the categorical act $U$ with trivial left $N$-action.

Let $X$ be the original collage space and $Y$ the core collage
on $U$. This gives an inclusion and retraction $Y\leftrightarrows X$.
The monoid $N$ has a right zero, so $BN$ is contractible. The
actual two-block pushout identifies $Y$, up to homotopy, with
the ordinary categorical act cone on $U$: its edge correspondence
is the product of the $S$-element space with $BN$, and projection
is an equivalence on every space of the pushout diagram.
Every positive rational product of $Y$ is therefore zero.

Both block spaces are contained in $Y$. The mixed-nerve
description makes the relative quotient $X/Y$ a suspension of
the quotient of their edge correspondence spaces. Suppose
$X\simeq\mathbb CP^2$. If $H^2(Y;\mathbb Q)=0$, all its
degree-two classes lift from that suspension and square to zero.
If $H^2(Y;\mathbb Q)\ne0$, the split injection from the
retraction identifies it with the one-dimensional $H^2(X;\mathbb Q)$,
whose generator again has zero square by pullback from $Y$.
Both cases contradict the projective-plane square.
\end{proof}

\begin{theorem}\label{C:seventeen-three-strong-component}
No three-object category with at most seventeen arrows and a
nontrivial strongly connected component can have classifying
space of type $\mathbb CP^2$.
\end{theorem}
\begin{proof}
A disconnected category is impossible. If an object is a retract
of another, deleting it removes at least three arrows and gives
a category with at most fourteen, contrary to
Theorem~\ref{C:CP2-lower17}. A strongly connected whole category
is excluded by Corollary~\ref{C:strong-category-CP2-minimum23}.
Thus there is a proper strong pair $S$ and one remaining monoid
corner $N$. Take opposites so that $S$ is the source block.
Write $s=|S|$, $m=|N|$ and $t=|F(a)|+|F(b)|$ for its two
nonempty cross stalks. Then $s\geq6$, $t\geq2$, and
$s+m+t\leq17$.

If $m=1$, apply Theorem~\ref{C:strong-pair-trivial-endpoint21}.
If both block spaces have zero second rational cohomology,
Lemma~\ref{C:two-block-zero-H2-products} excludes their square.

Suppose first that $H^2(BS;\mathbb Q)\ne0$.
The categorical Rees count gives $s\geq11$, so $m+t\leq6$.
For $m\geq3$, one has $m\leq4$ and $t\leq3$, and the
preceding lemma applies to the two stalks. If $m=2$, its
trivial right rational module is projective and $s\leq13$.
Nonzero second homology then forces a trivial-group $3\times3$
minimal band in $S$. Each object has at most two rows, and
Corollary~\ref{C:categorical-two-row-projective-sink} applies.

It remains that $H^2(BS;\mathbb Q)=0$ and
$H^2(BN;\mathbb Q)\ne0$. The latter requires $m\geq5$ by
the small-monoid classification, hence $s\leq10$.
If $t\leq3$, at least one cross stalk is a singleton.
Its state is fixed by $N$ and gives an actual section
$BN\to E$ of the edge-space projection $E\to BN$.
Since the other block has zero second cohomology, the whole
vertex degree-two group injects into that of $E$.
All degree-two classes of the pushout are therefore relative
to the vertex spaces and pull back from a suspension; their
products vanish.
We may consequently assume $t\geq4$, which forces $s\leq8$.

The minimal ideal of $S$ now has trivial group. If its total
row count is two, there is one row at each object and
Lemma~\ref{C:one-row-source-block-collapse} collapses $S$
to one trivial object. This is an actual homotopy-equivalent
ordinary left-act cone over $N$, with fewer than twenty-two
arrows, which is impossible.
The only remaining shape has total row and column counts
$(3,2)$ and $s=8$. Its row distribution is $(2,1)$ and its
column distribution is $(1,1)$; every nonidentity arrow of
$S$ is minimal. At its one-row object, every incoming internal
cross arrow is fixed by the minimal idempotent, and there are
no external incoming arrows. The one-sided deletion lemma
deletes that object from the entire collage. It removes at
least two corner arrows, three internal cross arrows, and
one external cross arrow. The resulting category has at most
eleven arrows, contrary to the global lower bound.
This exhausts all cases.
\end{proof}
